\chapter{Scheduler und Datenbeschaffung}\label{kap:scheduler}

Der Scheduler ist ein R-Prozess (\datei{RCode/updateScheduler.R}, Hauptschleife
in \datei{RCode/update_all_leagues_loop.R}). Er läuft dauerhaft im
Container und entscheidet, wann Daten geholt, simuliert und Seiten
geschrieben werden.

\section{Das Zeitfenster}

\begin{itemize}
  \item Aktiv ist der Scheduler täglich von \textbf{11:00 bis 23:00 Uhr}
        Berliner Zeit, also 720 Minuten. Der Beginn um 11 Uhr folgt aus dem
        Spielplan: 37~Prozent der Spiele der 2.~Frauen-Bundesliga beginnen
        vormittags um 11 Uhr.
  \item Die Laufzeit ergibt sich aus Fensterende minus Fensterbeginn. Die
        Umgebungsvariable \code{DURATION} bleibt leer. Ein fester Wert wäre
        eine zweite Quelle, die beim nächsten Verschieben des Fensters
        unbemerkt das Fenster beschneiden würde.
  \item Außerhalb des Fensters schläft der Prozess bis zum nächsten
        Fensterbeginn. Jeder Prozess ist ein Tageslauf: Er endet, sobald die
        nächste Runde nicht mehr vor 23 Uhr passt, und Docker startet den
        Container neu (\code{restart: unless-stopped}). Der neue Prozess
        schläft dann bis 11 Uhr -- pro Abend genau ein Neustart.
  \item Als „nach dem Fenster“ gilt deshalb schon ein Start ab
        \textbf{22:58 Uhr}, also wenn höchstens noch ein Takt (zwei Minuten)
        bleibt. Vorher plante ein Start um 22:58 bis 23:00 noch eine Runde,
        endete nach ihr und startete erneut -- rund sieben Neustarts je Abend,
        die \code{RestartCount} als Absturzanzeige unbrauchbar machten
        (Issue~\#287).
  \item Die Fensterrechnung erfolgt ausdrücklich in \code{Europe/Berlin},
        unabhängig von der Zeitzone des Prozesses.
\end{itemize}

\section{Takt und Live-Gating}

Im Normalbetrieb läuft eine Runde alle \textbf{120 Sekunden}, das sind bis
zu 361 Runden am Tag. Damit das API-Kontingent reicht, ist jede Runde zweistufig:

\begin{enumerate}
  \item \textbf{Live-Abfrage} (eine Anfrage für alle Ligen,
        \code{GET /v3/fixtures?live=78-79-80-…}): Läuft irgendwo ein Spiel?
  \item \textbf{Vollabruf} (eine Anfrage je Liga, also zehn) nur, wenn
        \begin{itemize}
          \item ein Spiel läuft, oder
          \item ein Spiel den Live-Feed verlassen hat, in den Saisondaten aber
                noch nicht als beendet steht (Issue~\#154), oder
          \item der letzte erfolgreiche Vollabruf länger als 3\,600 Sekunden
                zurückliegt (Sicherheitsnetz, Issue~\#226), oder
          \item die Live-Abfrage scheitert.
        \end{itemize}
\end{enumerate}

Die erste Runde eines Tages ist immer ein Vollabruf. Simuliert wird eine
Liga nur, wenn sich die Menge ihrer beendeten Spiele geändert hat. Verglichen
wird die Menge der Spiel-IDs, nicht ihre Anzahl. Seiten werden geschrieben,
wenn simuliert wurde oder sich etwas Sichtbares geändert hat, etwa ein
Zwischenstand oder eine Anstoßzeit.

\section{Das API-Kontingent}

Der Zugang läuft über den Pro-Plan von API-Football: 7\,500 Anfragen am Tag
und 300 je Minute. Ein ruhiger Tag verbraucht etwa 500 Anfragen, ein voller
Spieltag bis zu 3\,900 (Issue~\#190).

Nach jeder Anfrage liest der Scheduler die Kopfzeilen
\code{x-ratelimit-requests-remaining}, \code{-limit} und \code{-reset} mit.
Ein Regler (\datei{RCode/rate_limit_takt.R}) passt den Takt an das
Restkontingent an. Er hält den Abstand zwischen 60 und 5\,400 Sekunden,
plant mit einem Sicherheitsabschlag von 10~Prozent und schaltet mit 15~Prozent
Hysterese, damit der Takt nicht flattert. Bei gedrosseltem Takt zeigt die
Website einen Hinweis (Kapitel~\ref{kap:seiten}).

\section{Fehlertoleranz}

\begin{itemize}
  \item \textbf{Je Liga isoliert.} Scheitert der Abruf, die Aufbereitung oder
        die Simulation einer Liga, fällt nur diese Liga aus. Ihre letzte
        Prognose bleibt stehen, und alle anderen Seiten werden aktualisiert.
  \item \textbf{Zeitlimits.} 30 Sekunden für API-Football, 60 Sekunden für den
        Rust-Server, gegen hängende Verbindungen.
  \item \textbf{Startskript.} \datei{docker-start.sh} startet den Rust-Server,
        wartet auf dessen \code{/health}, startet ihn neu, falls er stirbt,
        und versucht den Scheduler bis zu fünfmal im Abstand von 30 Sekunden
        zu starten.
  \item \textbf{Warnungen sofort.} \code{options(warn = 1)} schreibt jede
        Warnung direkt ins Container-Log. Viele Fehler führen nicht zum
        Absturz, sondern zu einer fehlenden Sektion mit Warnung.
\end{itemize}

\section{Der Ablauf einer Runde}

\begin{enumerate}
  \item Kontingent prüfen, Takt festlegen.
  \item Live-Abfrage, bei Bedarf Vollabruf je Liga.
  \item Daten aufbereiten (\code{transform\_data}): Statusgruppen, Filter auf
        Hauptrunde, chronologische Sortierung.
  \item Für jede geänderte Liga \code{POST /simulate} mit \num{10000}
        Iterationen, in der 3.~Liga und der 2.~Frauen-Bundesliga zusätzlich
        der Malus-Lauf.
  \item Regionalliga-Kopplung und Aufstiegsspalten berechnen.
  \item Spieldetails je Liga über \code{POST /league-details}.
  \item \code{generate\_static\_site()} schreibt alle Seiten.
\end{enumerate}
